\part{Cadres d'action éthique et déontologique}
\chapter{La responsabilité en pratique : intervenir et codifier}
\label{chap:deontologie}

\section{Introduction : de la critique à l'action}
Les chapitres précédents ont dressé un constat critique : l'opacité des systèmes d'IA et leur incarnation d'une nouvelle forme de pouvoir bureaucratique posent des défis fondamentaux à la responsabilité. Ce chapitre, répondant aux deux dernières questions de recherche, passe de l'analyse à la proposition. Comment intervenir concrètement sur les algorithmes jugés prédateurs ? Et comment cultiver une culture de la responsabilité chez leurs concepteurs via un code de déontologie, en naviguant les «rapports de force» organisationnels ?

\section{Intervenir sur les algorithmes prédateurs : un cadre méthodologique}
Modifier un algorithme n'est pas un simple ajustement technique ; c'est un acte de gouvernance qui doit être éthiquement fondé et méthodologiquement rigoureux.

\subsection{Fondements éthiques : prudence et principes}
L'intervention se justifie par le «Principe Responsabilité» de \citet{Jonas1979} : face à l'incertitude des conséquences technologiques, la prudence impose de donner plus de poids à la «prophétie de malheur». La \textbf{Déclaration de Montréal pour une IA responsable} \citep{MontrealDeclaration2018} traduit cet impératif en dix principes actionnables (bien-être, autonomie, équité, etc.) qui serviront de boussole à notre démarche.

\subsection{Une démarche structurée en cinq étapes}
Nous proposons un processus itératif et participatif pour réaligner la logique d'un artefact technique avec des valeurs humaines.
\begin{enumerate}
    \item \textbf{Diagnostic et caractérisation du préjudice :} Identifier et qualifier précisément le préjudice (discrimination, manipulation, etc.) en mobilisant des méthodes mixtes (audits quantitatifs, entretiens qualitatifs).
    \item \textbf{Analyse causale via l'explicabilité (XAI) :} Utiliser les outils de l'XAI (LIME, SHAP) pour «ouvrir la boîte noire» et déterminer si le problème vient des données, de la fonction objectif, ou d'autres facteurs techniques.
    \item \textbf{Délibération éthique et redéfinition des objectifs :} Confronter les causes techniques aux principes de la Déclaration de Montréal. Cette étape, au cœur du processus, est politique et éthique : elle consiste à décider des compromis à faire (e.g., sacrifier un peu de «pertinence» pour plus de «diversité»).
    \item \textbf{Conception et implémentation de l'intervention :} Sur la base des nouveaux objectifs, implémenter les modifications techniques : correction des biais dans les données, ajout de contraintes d'équité à l'algorithme, ajout de mécanismes de contrôle pour l'utilisateur, etc.
    \item \textbf{Validation, déploiement et monitoring continu :} Tester rigoureusement la version modifiée, la déployer, et mettre en place une surveillance continue des métriques techniques et éthiques, ainsi que des canaux de feedback pour les utilisateurs.
\end{enumerate}

Ce processus transforme la modification d'un algorithme d'un «débogage» en un processus de gouvernance, réintroduisant la rationalité substantielle «les valeurs» dans un système dominé par la rationalité formelle «l'optimisation».

\section{Élaborer un code de déontologie pour les concepteurs}
Une approche plus proactive consiste à guider les pratiques en amont. Cependant, un code de déontologie doit reconnaître la réalité des «rapports de force» et ne pas faire peser toute la charge morale sur l'individu.

\subsection{Inspiration processuelle : le Biosphere Code}
L'inspiration principale du \textbf{Biosphere Code} \citep{BiosphereCode2015} n'est pas son contenu, mais son \textit{processus d'élaboration} : il a émergé d'un atelier collaboratif et interdisciplinaire.

C'est ce modèle qu'il faut suivre pour garantir la pertinence et la légitimité du code. Il ne s'agit pas d'imposer des règles d'en haut, mais de les co-construire avec toutes les parties prenantes (développeurs, designers, chefs de produit, juristes, éthiciens, usagers).

\subsection{Le code comme artefact culturel et outil de pouvoir}
Un code de déontologie n'est pas un document juridique. C'est un \textit{artefact culturel} dont le but est de :
\begin{itemize}[label=\textbullet]
    \item Créer une \textit{conscience professionnelle collective}.
    \item Fournir un \textit{langage commun} pour discuter des enjeux éthiques.
    \item \textit{Légitimer les préoccupations éthiques} au sein de l'organisation.
\end{itemize}
Il devient un outil pour les développeurs dans les «rapports de force» en leur permettant de passer d'une objection personnelle («je suis en désaccord») à un argument professionnel («ceci contrevient au principe 4 de notre code»). Il vise à combler le «accountability gap» non pas en désignant un coupable, mais en créant une culture de la \textit{responsabilité partagée}, comme le suggère l'esprit de l'article de \citet{Neuer2018}.

\subsection{Structure et principes directeurs proposés}
Le code pourrait s'articuler autour des principes suivants, ancrés dans les cadres internationaux :

\begin{table*}[ht]
\centering
\renewcommand{\tablename}{Tableau} 
\caption{Structure et principes d'un code de déontologie pour la personnalisation}
\label{tab:code_deonto}
\begin{tabular}{@{}p{0.25\textwidth}p{0.65\textwidth}@{}}
\toprule
\textbf{Principe directeur} & \textbf{Devoirs et implications pratiques pour le concepteur} \\ \midrule
\textbf{1. Primauté du bien-être et non-malveillance} & Évaluer les risques d'addiction, d'anxiété ou d'exploitation. Refuser et dénoncer les «dark patterns». Privilégier le bien-être à long terme sur l'engagement à court terme. \\
\textbf{2. Équité et non-discrimination} & Auditer systématiquement les données et les modèles pour détecter et atténuer les biais. Utiliser des métriques d'équité adaptées. \\
\textbf{3. Transparence et explicabilité} & Documenter rigoureusement les choix de conception, les données et les limites. Utiliser des techniques XAI. Contribuer à fournir des informations claires aux utilisateurs. \\
\textbf{4. Respect de l'autonomie et vie privée} & Appliquer les principes de minimisation des données et de «Privacy by Design». Offrir aux utilisateurs un contrôle réel et simple sur leurs données et les paramètres. \\
\textbf{5. Responsabilité et redevabilité} & Assumer la responsabilité de la qualité et des impacts des algorithmes. Coopérer aux audits. Exercer un devoir d'alerte face aux risques éthiques. \\
\textbf{6. Supervision humaine} & Concevoir les systèmes pour permettre une supervision et une intervention humaines appropriées, en particulier pour les décisions à fort impact. \\ \bottomrule
\end{tabular}
\end{table*}

\section{Conclusion du chapitre}
Face à la puissance de la bureaucratie algorithmique, l'inaction n'est pas une option. Ce chapitre a esquissé deux voies d'action complémentaires. La première est une méthodologie d'intervention réactive, structurée et délibérative, pour corriger les systèmes prédateurs existants. La seconde est une approche proactive et culturelle : la co-construction d'un code de déontologie qui arme les professionnels pour intégrer l'éthique au cœur de leur pratique et pour naviguer les pressions organisationnelles. Ensemble, ces approches visent à passer d'une culture de l'optimisation à tout prix à une culture de la responsabilité partagée.
